\chapter{Petunjuk dan Kriteria Penilaian Latihan}
\label{app:petunjuk}

Kunci jawaban ada di Lampiran~D, tetapi bentuknya sengaja tidak lengkap: angka akhir
dan sketsa langkah, bukan pembuktian penuh. Alasannya konsistensi. Angka akhir tanpa
jejak pemeriksaan akan mengajari pembaca kebiasaan yang justru ingin dihapus buku ini,
dan bukti yang disodorkan utuh menghapus satu-satunya latihan yang berguna, yaitu
mencari sendiri di mana argumen belum menutup. Jadi pakai kunci ini sebagai pembanding
terakhir, bukan sebagai tempat memulai. Yang disediakan di sini adalah cara memeriksa
pekerjaan sendiri sebelum kamu membuka Lampiran~D.

\section*{Tiga jenis latihan dan cara memeriksanya}

\begin{center}
\small
\begin{tabular}{@{}p{2.6cm}p{5.6cm}p{5.8cm}@{}}
\toprule
\textbf{Jenis} & \textbf{Yang dinilai} & \textbf{Cara memeriksa} \\
\midrule
Bukti & setiap langkah harus berangkat dari definisi atau hasil yang sudah
  dibuktikan; tidak boleh ada lompatan & tukar naskah dengan teman; minta dia
  menunjuk kalimat yang tidak bisa ia turunkan sendiri \\
Hitung & angka akhir dan langkah antara & ulangi dengan Python; hitung galat
  relatif terhadap hasil eksak; periksa apakah jawaban masuk akal secara
  orde besaran \\
Nilai klaim & mutu argumen, bukan benar-salah binari & tuliskan lebih dulu apa
  yang harus benar agar klaim itu berlaku; cari satu contoh yang mematahkannya;
  baru putuskan "benar", "salah", atau "belum cukup bukti" \\
\bottomrule
\end{tabular}
\end{center}

\section*{Alat verifikasi per bagian}

\begin{center}
\small
\begin{tabular}{@{}llp{7.2cm}@{}}
\toprule
\textbf{Bab} & \textbf{Alat} & \textbf{Pemeriksaan khas} \\
\midrule
1--5 & \texttt{python3} & tabel nilai, faktorisasi, jumlah parsial; bandingkan
  dengan taksiran analitik (mis. $H_n\approx\ln n+\gamma$) \\
6--10 & \texttt{numpy}, \texttt{scipy.linalg} & rank, ruang nol (\texttt{null\_space}),
  nilai eigen, SVD, kuadrat terkecil (\texttt{lstsq}); periksa
  $\|A\mathbf{v}-\lambda\mathbf{v}\|$ \\
11--16 & \texttt{scipy.integrate}, beda hingga & kuadratur vs antiturunan; laju
  galat pada beberapa $h$ atau beberapa $n$ \\
17--21 & \texttt{numpy.random}, \texttt{scipy.stats} & simulasi Monte Carlo dengan
  \emph{seed} tetap; bandingkan peluang empiris dengan peluang eksak dan
  laporkan galat bakunya \\
22 & \texttt{struct}, \texttt{float.hex} & bitfield, ULP, epsilon mesin;
  bandingkan algoritme naif dengan Kahan \\
23--24 & \texttt{scipy.optimize} & \texttt{linprog}, \texttt{minimize},
  \texttt{milp}; periksa syarat KKT dan dualitas kuat secara numerik \\
25--26 & \texttt{numpy} & pemeriksaan gradien dengan beda hingga; kurva galat
  latih vs validasi pada beberapa setelan regularisasi \\
27--30 & \texttt{numpy}, \texttt{scikit-learn} bila ada & matriks konfusi dari
  prediksi tersimpan; interval Wilson untuk proporsi; akurasi di bawah serangan
  untuk beberapa $\varepsilon$ dan beberapa \emph{seed} \\
\bottomrule
\end{tabular}
\end{center}

\section*{Cara menilai pekerjaan sendiri}

\begin{enumerate}
  \item Tuliskan lebih dulu apa yang diketahui dan apa yang ingin dibuktikan.
        Bila keduanya belum jelas, kesulitannya bukan pada matematika.
  \item Untuk latihan hitung, tebak jawaban sebelum menghitung. Tebakan yang salah
        tetap berguna: ia menunjukkan di mana intuisi meleset.
  \item Bandingkan hasil dengan dua cara berbeda. Satu cara bisa benar karena
        kebetulan; dua cara yang sepakat jarang kebetulan.
  \item Bila hasilnya berbeda jauh, cari kesalahan pada kasus sederhana lebih dulu
        (satu variabel, satu langkah, $n=1$).
  \item Tuliskan batas keberlakuan jawaban. Jawaban benar untuk kasus khusus yang
        dinyatakan sebagai jawaban umum adalah kesalahan yang paling sering lolos.
\end{enumerate}
